\begin{comment}
\subsubsection{Vague quantifiers as indirect indicators of group-specific norms}
In the 2002 and 2004 NSS response categories for the outcome variables are coded using qualitative frequency categories such as ``very often,'' ``not very often,'' etc. rather than more ``objective'' numerical estimates (``once a week'', ``two to four times a month,'' etc.)  These types of survey items are what is known in the questionnaire design literature as ``vague quantifiers'' \citep[74]{schaeffer_presser03}. The use of qualitative frequency estimates is generally thought of as problematic by researchers who want to take these response variables as indicators of some underlying ``true'' value or frequency.  The reason why vague quantifiers are problematic for the particular purpose of using them as indices of objective   behavior frequencies should be obvious:  if different groups of people attach different underlying numeric frequencies to the same qualitative label, then this introduces strict non-comparability across persons. For instance, one person's idea of the difference between ``very often'' and ``often'' may differ systematically from another person depending on their value on a predictor variable.\footnote{The only qualitative frequency indicator that appears to have an unambiguous meaning across respondents is, as should be expected, ``never.''} 

But this is not the only way to approach the issue of vague-quantifiers.  In what follows, I take advantage of recent research in cognitive and social psychology that reveals a key substantive property of individual ``biases'' in connecting numeric frequencies to qualitative labels.  This research reveals that under a fairly general set of conditions, group heterogeneity in the interpretation of qualitative frequency categories is connected to the presence of structured variation in group-specific \emph{standards} of what is the general norm, standards generated by both biases in the base rate at which persons in those groups engage in the activity in question and in their underlying perception as to how often they should engage in that activity. Persons use the modal behavior of their peer-group as a \emph{cognitive anchor} \citep{oliver80} to produce a judgment of what numeric frequency of a given practice counts as ``often'' or ``not very often''; if the behavior is rare within the peer-group, then lower numeric frequencies may be judged as worthy of the label ``very often'' than if the behavior is statistically normative within the peer-group. This interpretation sheds light on socio-demographic differences keyed to gender, class, race and education in the interpretation of subjective frequency estimates across a wide range of outcomes, include culture and leisure consumption activities (\citealp[481]{wright_etal94}; \citealp[406]{schaeffer91}; \citealp[304-305]{wanke02}) Thus, heterogeneity in the interpretation of vague quantifiers can most profitably be thought of as an (indirect) window into the \emph{culture talk norms} characteristic of different status cultures.\footnote{\citet{wright_etal94} show that, in the case of frequency of television watching, the assignment of numerical values to qualitative categories differs across social class, with those class fractions more likely to watch television assigning higher numeric frequencies to the same qualitative category (e.g. ``often'') than those who did not watch a lot of television.   This is an observation that has been replicated by \citet{wanke02} in the case of culture consumption behavior: socio-demographic groups who have a higher base propensity to engage in the activity more frequently (e.g. the young or the higher-educated) have a ``tougher'' (shifted to the left of the distribution of counts) threshold for what constitutes ``very often'' than those groups who are less likely to engage in the activity (as given by the self-reported counts). }

\subsubsection{Generalized Ordered Probit Model with Cutpoint-Shift}
Following this line of reasoning, I conceptualize the individual's qualitative response to each of the culture talk items as the observed realization of the person's unobserved ``objective'' (e.g. characterizable in numeric terms) frequency of culture talk as filtered through the individual's propensity to assign his or her own numeric range to specific qualitative categories available in the survey.  I assume that each person in the sample has an ascertainable underlying propensity to engage in culture talk, but that persons in different socio-demographic categories partition this continuum (e.g. assign specific ``cutpoints'' to the distribution) in systematically different ways.  I do not model this variation as either random or as unique to each individual.  Instead, I hypothesize that variation in the extent to which persons assign distinct numeric frequencies to qualitative categories is structured by socio-economic status.  I use a generalization of the ordered probit model that allows for the direct estimation of socio-demographic variation in cutpoint assignment in order to investigate there are status-linked differences in the perceived normativity of culture talk.  

Like the standard ordered probit model, the generalized ordered probit model presumes that the observed qualitative response is the realization of an underlying continuous variable ($Y_{it}^*$)---representing the unobserved objective frequency----that is segmented at (in this case) $M-1=3$ cutpoints ($\alpha_{jt}$):

\begin{equation}
	Y_{it}=g(Y_{it}^{*})=\left\{\begin{array}{rl}
	  1 &\mbox{if $Y_{it}^{*} \leq \alpha_{1t}$} \\
	  2 &\mbox{if $\alpha_{1t} < Y_{it}^{*} \leq \alpha_{2t}$}\\ 
	  3 &\mbox{if $\alpha_{2t} <Y_{it}^{*} \leq \alpha_{3t}$}\\ 
	  4 &\mbox{if $\alpha_{3t} \leq Y_{it}^{*}$} 
	 \end{array} \right.
	 \label{eq:equation3}
\end{equation}
 
We only observe $Y_{it}=j$ if $\alpha_{jt-1} < Y_{it}^{*} \leq \alpha_{j}$.  In the traditional ordered probit model, the $\alpha_{jt}$ are a constant across all respondents.  That is, every person is assumed to map their underlying frequencies to the same qualitative responses. We  have seen that this assumption is unlikely to hold for these data. The generalized ordered probit model \citep{terza85} relaxes this restriction to account for group-specific biases in the interpretation of the qualitative frequency categories. This model allows for $M-2$ of the cutpoints (where $M$ is the total number of ordered categories) to be a function of a selected set of socio-demographic variables:

\begin{equation}
\alpha_{jt}=\hat{\alpha_{jt}}+\theta_{jt} E_i+\sum_k \phi_{jkt}Z_{k}
\end{equation}•

In the cutpoint equations,  the sign of the coefficient estimate tells us whether persons who are coded as high in the value of the $k^{th}$ variable tend to shift the value of their $j^{th}$ cutpoint to the right or the left of the underlying distribution for a given culture talk domain $t$. The substantive interpretation of the coefficients corresponding to each predictor variable in the cutpoint equations is straightforward.  A \emph{positive} coefficient estimate indicates that persons with high values on that predictor tend to shift that cutpoint to the \emph{right} of the underlying distribution; that is they apply more stringent standards to similar underlying quantitative frequencies (for purposes of categorizing them, for instance, as ``very often'') than persons with low values in that predictor.  A \emph{negative} coefficient indicates the opposite: for individuals with high value on this predictor, the cutpoints are shifted to the \emph{left}; these persons judge lower underlying frequencies (as sufficient for categorization as, for instance,``fairly often'') than persons with low values on this predictor.    Following this line of reasoning, we should expect right-ward shift in cutpoint assignment among high-status elites; that is, $\theta_{jt}>0$.
\end{comment}